\subsubsection*{Argument Forms as a Bridge to Formal Systems}

\begin{remark*}[Working vocabulary]
This topic distinguishes semantic validity from formal derivability and explains
why formal proof systems are deferred.
\end{remark*}

\begin{remark*}[Exposition]
A valid argument form is a semantic object: it preserves truth under all truth
assignments. A formal rule of inference is a syntactic object inside a specified
proof system. Many formal systems include rules corresponding to the valid
argument forms justified in this section, but the semantic justification and the
formal rule are not the same object.
\end{remark*}

\begin{remark*}[Exposition]
Formal proof systems require additional machinery: formal derivations,
assumptions, discharge rules, axioms or axiom schemata, and precise proof-checking
rules. Natural deduction, Hilbert systems, tableaux, resolution, and sequent
calculus are examples of such systems. They are important, but they are not the
main purpose of this introductory propositional-logic section.
\end{remark*}

\begin{remark*}[Transition]
The tools validated here are safe to use in ordinary mathematical proof because
they are truth-preserving. Later, when formal proof systems are introduced, these
same patterns can be represented as formal rules, derived rules, or admissible
moves depending on the chosen system.
\end{remark*}
